% !PRINTED 216
\weberpage{216}
%<<CONT
in it $\gamma, \delta$ can be replaced by $\gamma + h\alpha$, $\delta + h\beta$,
and by letting $h$ run through a complete system of residues modulo $n$,
one recognises that to every pair of numbers $\alpha, \beta$ there belong $n$
pairs of numbers $\gamma, \delta$ satisfying condition (23). Accordingly the
% [sic] The German source prints (23); the determinant condition to which
% the pairs gamma, delta must conform is (22).
order of the group $\mathfrak{B}$ is:
\[
\nu = n^3 \Pi\left(1 - \frac{1}{p^2}\right),
\]
or, if we further introduce the numerical function

\weq{25}{\psi(n) = n\Pi\left(1 + \frac{1}{p}\right)}
then:

\weq{26}{\nu = n\varphi(n)\psi(n)}
and the order of the group $\mathfrak{A}$:

\weq{27}{\mu = n\varphi(n)^2\psi(n).}

\begin{center}
\S\,64. \gesperrt{The irreducible factors of the division equation.}
\end{center}
If
\[
\begin{aligned}
\nu &= \partial\mu - b\mu'\\
\nu' &= -c\mu + a\mu'
\end{aligned}
\]
and $a\partial - bc$ is relatively prime to $n$, then the greatest common
divisor of $\mu, \mu', n$ is at the same time the greatest common
divisor of $\nu, \nu', n$. The roots $x_{\mu,\mu'}$ of the division equation
therefore split, according to the greatest common divisor of
$\mu, \mu', n$, into systems which under the substitutions of the group $\mathfrak{A}$
always pass only into one another, that is, the group $\mathfrak{G}$ is \gesperrt{intransitive}, % Source verified: the group letter is Fraktur G.
and the systems of \gesperrt{intransitivity} are the
systems $x_{\mu,\mu'}$ in which $\mu, \mu', n$ have one and the same greatest
common divisor. The division equation is therefore
reducible (except when $n$ is a prime) and the roots
of one of these systems of intransitivity % [source terminology] German switches here from systems of intransitivity to imprimitivity.
satisfy a rational
equation (Bd.\,I, \S\,157).

According to the preceding section there are $\varphi(n)\psi(n)$
pairs of numbers $\mu, \mu'$ whose greatest common divisor with $n$
is equal to 1, and the roots corresponding to these pairs of numbers
$x_{\mu,\mu'}$ therefore satisfy a rational equation of degree $\varphi(n)\psi(n)$.
This equation we call the \gesperrt{proper division equation}
for the divisor $n$, because only when $\mu, \mu', n$ are without a common
divisor is $x_{\mu,\mu'}$ not at the same time a root of a division
%CONT>>
